\section{The non-normal fibre: exact claim under investigation}

\paragraph{Status convention.}
This notebook does not assert that the CDP theorem is false. It studies one
specific inference in the printed proof and the conditional application of a
conductor-supported differential to the candidate fibre $W$. A defect in the
printed inference, an alternative proof of the published theorem, and failure
of the candidate construction are logically distinct possibilities.

\subsection{The local conductor proposition}

\begin{proposition}[conductor-supported ambient differential]
\Conditional\ Let $g:\mathcal X\to\Delta$ be proper, with $\mathcal X$ a
smooth complex threefold, and let $S=g^{-1}(0)$ be a reduced
normal-crossings surface. Write $\eta:\widetilde S\to S$ for the
normalization, $D\subset S$ for the reduced singular locus, and
$\widetilde D=\eta^{-1}(D)$. Let $A$ be a line bundle on $S$ such that
$\eta^*A\simeq\mathcal O_{\widetilde S}$, and choose a nowhere-vanishing
section $e$ of $\eta^*A$. Then there is a nonzero section
\[
 s\in H^0\!\left(S,\Omega^1_{\mathcal X}|_S\otimes A\right)
\]
whose image in $\Omega^1_S\otimes A$ is supported on $D$, hence is
torsion, and whose image in
\[
 \widetilde\Omega^1_S\otimes A,
 \qquad
 \widetilde\Omega^1_S:=\Omega^1_S/\mathcal T(\Omega^1_S),
\]
is zero. If $S$ is connected and $A\not\simeq\mathcal O_S$, then
$H^0(S,N^*_{S/\mathcal X}\otimes A)=H^0(S,A)=0$, so $s$ does not come
from the conormal term.
\end{proposition}

\begin{proof}
Locally $g=ux$, $uxy$, or $uxyz$, where $u$ is a unit. The conormal map is
generated by $dg|_S$. At a double point,
\[
 dg|_S=u(y\,dx+x\,dy),
\]
and at a triple point it has coefficients in $(xy,xz,yz)$. Thus $dg|_S$
is nonzero away from $D$ but restricts to zero on $D$; equivalently, it
lies in the conductor ideal times $\Omega^1_{\mathcal X}|_S$.

The normalization--conductor square gives
\[
0\longrightarrow\mathcal O_S\longrightarrow
 \eta_*\mathcal O_{\widetilde S}\oplus\mathcal O_D
 \longrightarrow\eta_*\mathcal O_{\widetilde D}\longrightarrow0.
\]
After tensoring with the locally free sheaf
$E=\Omega^1_{\mathcal X}|_S\otimes A$, the compatible pair
\[
 \bigl(\eta^*(dg|_S)\otimes e,\,0\bigr)
\]
descends uniquely to $s$. It is nonzero off $D$. Its intrinsic differential
vanishes off $D$, so its image in $\Omega^1_S\otimes A$ is
conductor-supported torsion. For nontrivial $A$, a global section of $A$
pulls back to constants on the compact connected components of
$\widetilde S$. Compatibility across the conductor and connectedness of
$S$ force either all constants to vanish or all to be nonzero; the latter
would trivialize $A$. Hence $H^0(S,A)=0$.
\end{proof}

Consequently the implication
\[
\left[
 H^0(S,N^*_{S/\mathcal X}\otimes A)=0,\quad
 H^0(S,\widetilde\Omega^1_S\otimes A)=0
\right]
\Longrightarrow
H^0(S,\Omega^1_{\mathcal X}|_S\otimes A)=0
\]
is not a formal consequence of the conormal sequence for a non-normal
normal-crossings surface. The official CDP article prints precisely this
reduction on p.~692. This observation concerns that inference, not by itself
the truth value of the theorem.

\subsection{Conditional application to the candidate fibre}

\Conditional\ Assume the following upstream assertions about the candidate:
\begin{enumerate}[label=(\roman*)]
\item $f:X\to B$ is a proper nonconstant holomorphic map from a smooth
threefold to a smooth curve;
\item $W=f^{-1}(p_0)$ is the reduced scheme-theoretic fibre and is a
normal-crossings Cartier divisor;
\item $\eta:\widetilde W\to W$ has
$\widetilde W\simeq\mathrm{dP}_6$;
\item $H^2(X;\mathbb Z)=0$.
\end{enumerate}

For any $L\in\operatorname{Pic}(X)$, put
\[
 \mathcal F_L=T_X\otimes L,
 \qquad
 A_L=(L^\vee\otimes\omega_X)|_W.
\]
Because $H^2(X;\mathbb Z)=0$, $c_1(L^\vee\otimes\omega_X)=0$.
Since $H^1(\mathrm{dP}_6,\mathcal O)=0$, the Chern-class map
$\operatorname{Pic}(\mathrm{dP}_6)\to H^2(\mathrm{dP}_6;\mathbb Z)$
is injective. Therefore
\[
 \eta^*A_L\simeq\mathcal O_{\widetilde W}.
\]
The conductor proposition yields
\[
 H^0(W,\Omega_X^1|_W\otimes A_L)\ne0.
\]

\subsection{The exact duality implication required}

As an effective Cartier divisor in a smooth threefold, $W$ is a pure
two-dimensional Cohen--Macaulay Gorenstein complex space. The resolution
\[
0\longrightarrow\mathcal O_X(-W)\longrightarrow\mathcal O_X
 \longrightarrow\mathcal O_W\longrightarrow0
\]
gives
\[
 \omega_W\simeq(\omega_X\otimes\mathcal O_X(W))|_W.
\]
Because $W$ is a full fibre,
$\mathcal O_X(W)|_W\simeq\mathcal O_W$, hence
$\omega_W\simeq\omega_X|_W$. Absolute analytic
Serre--Grothendieck duality must then specialize to the perfect pairing
\[
 H^2(W,\mathcal F_L|_W)^\vee
 \simeq
 \operatorname{Hom}_W(\mathcal F_L|_W,\omega_W)
 =
 H^0(W,\Omega_X^1|_W\otimes A_L).
\]
Thus the conductor section implies
\[
 H^2(W,\mathcal F_L|_W)\ne0.
\]
Normality is not required here; compactness, Cohen--Macaulayness,
Gorenstein dimension $2$, and local freeness of $\mathcal F_L|_W$ are the
relevant hypotheses.

\begin{obligation}[duality specialization]\label{obl:duality}
\Open\ Verify in the chosen analytic duality reference that the compact
possibly non-normal space $W$ has the stated dualizing sheaf and trace, and
that the resulting pairing is perfect for the locally free sheaf
$\mathcal F_L|_W$. Check the Cartier resolution, adjunction, the triviality
of $\mathcal O_X(W)|_W$, the shift by dimension $2$, and the final
identification
\[
(\mathcal F_L|_W)^\vee\otimes\omega_W
 \simeq \Omega_X^1|_W\otimes A_L.
\]
This obligation is the exact bridge from the conductor section to
$H^2(W,\mathcal F_L|_W)\ne0$.
\end{obligation}

\subsection{The exact top-degree base-change implication required}

The map $f$ is flat if a local parameter on $B$ pulls back to a nonzero
divisor in every regular local ring $\mathcal O_{X,x}$. The vector bundle
$\mathcal F_L$ is then $f$-flat. Since the fibres have complex dimension
$2$,
\[
 H^3(X_c,\mathcal F_L|_{X_c})=0
\]
for every full scheme-theoretic fibre $X_c$. Applying the analytic
cohomology-and-base-change theorem in degree $q=3$ gives base change in
degree $2$; more precisely,
\[
 \beta^2_c:
 (R^2f_*\mathcal F_L)_c\otimes_{\mathcal O_{B,c}}\mathbb C(c)
 \xrightarrow{\sim}
 H^2(X_c,\mathcal F_L|_{X_c}).
\]
This is the precise consequence of
B\u{a}nic\u{a}--St\u{a}n\u{a}\c{s}il\u{a}, Chapter~III,
Section~3, Theorem~3.4 and Corollary~3.5, that should replace the
manuscript's broad citation. At $c=p_0$, the nonzero target implies
\[
 (R^2f_*\mathcal F_L)_{p_0}\ne0,
 \qquad\text{hence}\qquad
 R^2f_*(T_X\otimes L)\ne0.
\]
Only surjectivity of $\beta^2_{p_0}$ is needed for this nonvanishing; the
isomorphism is needed for the later identification of the support with the
fibrewise $H^2$-jumping locus.

\begin{obligation}[top-degree base change]\label{obl:base-change}
\Open\ Prove properness and flatness near $p_0$, verify that
$\mathcal F_L$ is $f$-flat and that $W$ is the full scheme-theoretic
fibre, and state the cited analytic theorem in the form that turns the
vanishing of fibrewise $H^3$ into the displayed isomorphism $\beta^2_c$.
At minimum, prove surjectivity of $\beta^2_{p_0}$; do not replace this with
semicontinuity alone.
\end{obligation}

\subsection{Present evidentiary status}

\begin{itemize}
\item The local conductor calculation was independently checked and called
sound by Reviewers A and C.
\item Reviewer A found the application plausible but requested exact
duality and base-change hypotheses and citations.
\item Reviewer B classified the corresponding lemma and theorem as
secondary and reviewed them only for consistency; Reviewer B did not
independently certify the analytic bridge.
\item Reviewer C called the application coherent and the printed-proof
issue substantiated, while explicitly separating it from falsity of the
published theorem and from validity of the candidate construction.
\item The application to $W$ remains conditional on the upstream
construction of $X$, the asserted fibre geometry, and
$H^2(X;\mathbb Z)=0$, none of which the three audits collectively certify.
\end{itemize}

\subsection{Verification and falsification programme}

\begin{enumerate}[label=\arabic*.]
\item \textbf{Local algebra.} Work in
$\mathbb C\{x,y,z\}/(xy)$ and $\mathbb C\{x,y,z\}/(xyz)$. Verify the
conductor ideals, injectivity of
$\mathcal O_S\,dg\to\Omega^1_{\mathcal X}|_S$, exactness of the
normalization--conductor sequence at double and triple points, and the
non-conormal character of $s$.

\item \textbf{Independent reconstruction of $W$.} Verify directly from
the quotient charts that the full fibre is reduced, Cartier and
normal-crossings, that its normalization is exactly $\mathrm{dP}_6$, and
that the stated conductor curves and triple points are exhaustive.

\item \textbf{Picard calculation.} Recompute
$\eta^*((L^\vee\otimes\omega_X)|_W)$ from explicit transition functions,
rather than only from the claimed global topology. A single $L$ for which
this pullback is nontrivial would falsify the universal application of the
conductor proposition.

\item \textbf{Duality audit.} Derive the dualizing complex of $W$ from
the Cartier resolution, check the shift and trace pairing, and independently
identify $(T_X\otimes L)|_W^\vee\otimes\omega_W$. Failure of
Cohen--Macaulayness, Cartierness, compactness, adjunction, or the pairing
would break the claimed implication.

\item \textbf{Base-change audit.} Prove properness and flatness near
$p_0$, verify that $\mathcal F_L$ is $f$-flat and that the fibre used is
the full scheme-theoretic fibre, then apply the stated degree-$3$-to-
degree-$2$ theorem. Alternatively, construct the local finite-free complex
representing $Rf_*\mathcal F_L$ and compute $\beta^2_{p_0}$ directly.

\item \textbf{Literature comparison.} Reproduce the CDP p.~692 argument
with exactly the same fibre, twist, and definition of
$\widetilde\Omega^1_S$. Determine whether an omitted hypothesis restores
the inference or whether an independent proof bypasses it. Either outcome
must be recorded without converting a proof-gap diagnosis into a theorem
verdict.
\end{enumerate}

\begin{obligation}[CDP comparison]\label{obl:cdp-comparison}
\Open\ Produce a line-by-line diagram comparing the candidate's
ambient, conormal, and intrinsic cotangent sequences with the exact sheaves
in the cited CDP reduction. Record separately: the validity of the printed
inference, the possibility of an independent repair, and the validity of
the candidate construction.
\end{obligation}

\begin{test}[decisive outcomes]
The Section~10 route fails if the descended section $s$ is zero after the
true conductor identifications, if it lies in the conormal image, if the
duality specialization is unavailable, or if fibrewise nonvanishing does
not produce the asserted nonzero stalk of $R^2f_*(T_X\otimes L)$.
Conversely, proving \cref{obl:duality,obl:base-change,obl:cdp-comparison}
would remove the printed CDP reduction as an automatic contradiction, but
would still leave the construction and topology obligations of
\cref{obl:finite-model,obl:cusp-collapse,obl:vankampen,obl:leray}.
\end{test}
